\documentclass[10pt,a4paper]{amsart}
\usepackage[margin=19mm]{geometry}
\usepackage{amsmath,amssymb,mathtools}
\usepackage[scr=rsfs,cal=euler]{mathalfa}
\usepackage{xfrac}
\usepackage[unicode,hidelinks,bookmarks=false]{hyperref}
\newcommand{\ie}{i.e.\ }
\newcommand{\cf}{cf.\ }
\newcommand{\resp}{respectively\ }
\newcommand{\Subsection}{\subsection}
\providecommand{\nobreakdash}{\nobreak\hbox{-}}
\setlength{\emergencystretch}{2em}
\multlinegap0pt
\usepackage{amssymb}
\usepackage{bm}
\usepackage{mathtools}
\usepackage{float}
\usepackage{algorithm}\usepackage{algpseudocode}
\newtheorem{lemma}{Lemma}
\newtheorem{proposition}{Proposition}
\title{Urban transit network design using spanning tree: A case study of Canberra transit network}
\author{Satoshi Suguira, Kam-Fung Cheung, Michael G. H. Bell, Hitomi Nakanishi, Fumitaka Kurauchi, Supun Perera, Yogi Vidyattama}
\date{Source: arXiv:2604.15754v1 (CC BY 4.0)}
\begin{document}
\maketitle

\noindent\emph{Source and licence.} Urban transit network design using spanning tree: A case study of Canberra transit network. Version arXiv:2604.15754v1 (CC BY 4.0), distributed by arXiv under the Creative Commons Attribution 4.0 International licence, http://creativecommons.org/licenses/by/4.0/, as recorded in the arXiv licence metadata for this exact version and shown on the abstract page https://arxiv.org/abs/2604.15754v1. Full licence notice: \texttt{licenses/arxiv-2604.15754-CC-BY-4.0.txt}.\par\smallskip

\noindent\textit{Editorial scope.} Counted unit: the article's proposition prop:swapping\_procedure (source0.tex lines 339--342). The article's complete original proof of it is delivered with it (source0.tex lines 343--345, one \texttt{proof} environment of 3 lines). No mathematics is written, repaired or re-derived: the statement and the proof are the article's own lines, in the article's own notation, and no retained line is substituted or re-flowed.  Retained uncounted as the source-local context the proof needs: lines 307--310; lines 331--334.  Nothing was omitted from any retained span, so the package carries no omission record.  No retained line contains a citation, so the wrapper carries no bibliography and no bibliography entry.  No retained line contains a figure environment, an image-inclusion command or drawing code, so no image is delivered and the package declares no asset dependency.  Editorial formatting only, in addition: an AMS \texttt{amsart} preamble that restates the article's own declarations and macros used by the retained lines, namely the environments \texttt{lemma}, \texttt{proposition} on separate counters, together with the standard packages those lines need.  The retained environments print in this document as Lemma 1, Proposition 1 in order of appearance, whereas the article prints the same items with the numbering of its own full document.  The complete native source of the article as distributed in the arXiv e-print for this exact version is bundled unchanged as \texttt{sources/198641/source0.tex}.
\begin{lemma}
  A spanning tree $T\left(V,E\right)$ is always divided into two disjoint subnetworks when a link is deleted from it, where $V$ and $E$ are set of nodes and links in the spanning tree, respectively.
  \label{lem:tree_divided_two_subnetworks}
\end{lemma}

\begin{lemma}
  The graph obtained by connecting one node in the spanning tree ${^1}T_{-a}$  and another node in the spanning tree ${^2}T_{-a}$  is a spanning tree, where ${^1}T_{-a}$  and ${^2}T_{-a}$  are trees by removing link $a$ in the spanning tree $T\left(V,E\right)$.
  \label{lem:connected_graph_is_tree}
\end{lemma}

\begin{proposition}
  A spanning tree is always partitioned into two subnetworks when a link is removed from it, and the graph obtained by adding a new link that connects a pair of nodes with one in each subnetwork is a spanning tree.
  \label{prop:swapping_procedure}
\end{proposition}

\begin{proof}
  A direct result from Lemma~\ref{lem:tree_divided_two_subnetworks} and Lemma~\ref{lem:connected_graph_is_tree}.
\end{proof}

\end{document}
